\documentclass[12pt,a4paper,austrian]{article}
\usepackage{graphicx}
\usepackage{booktabs}
\usepackage{float}
\usepackage[austrian, english]{babel}
\selectlanguage{english}
\usepackage[utf8]{inputenc}
\usepackage{listings}
\usepackage{multirow}
\usepackage{epstopdf}
\usepackage{amsmath}
\usepackage{amssymb}
\graphicspath{{./fig/}}

%% Satzspiegel
\setlength{\hoffset}{-1in} \setlength{\textwidth}{18cm}
\setlength{\oddsidemargin}{1.5cm}
\setlength{\evensidemargin}{1.5cm}
\setlength{\marginparsep}{0.7em}
\setlength{\marginparwidth}{0.5cm}

\setlength{\voffset}{-1.9in}
\setlength{\headheight}{12pt}
\setlength{\topmargin}{2.6cm}
   \addtolength{\topmargin}{-\headheight}
\setlength{\headsep}{3.5cm}
   \addtolength{\headsep}{-\topmargin}
   \addtolength{\headsep}{-\headheight}
\setlength{\textheight}{27cm}

\setlength{\floatsep}{12 pt plus 0 pt minus 8 pt}
\setlength{\textfloatsep}{12 pt plus 0pt minus 8 pt}
\setlength{\intextsep}{12 pt plus 0pt minus 8 pt}

\tolerance2000
\emergencystretch20pt

\newcommand{\eg}[1]{\selectlanguage{english}\textit{#1}\selectlanguage{english}}
\newcommand{\filename}[1]{\begin{small}\texttt{#1}\end{small}}

\newcommand\IFT{\unitlength1mm\begin{picture}(10,2) \put (1,1)
{\circle{1.7}} \put(2,1){\line(1,0){5}} \put(8,1)
{\circle*{1.7}}\end{picture}}
\newcommand\FT{\unitlength1mm\begin{picture}(10,2) \put (1,1)
{\circle*{1.7}} \put(2,1){\line(1,0){5}} \put(8,1)
{\circle{1.7}}\end{picture}}

\newcommand{\choicebox}{\fbox{\rule{0pt}{0.5ex}\rule{0.5ex}{0pt}}}

\newenvironment{wahrfalsch}%
  {\bigskip\par\noindent\makebox[1cm][c]{true}\hspace{3mm}\makebox[1cm][c]{false}
   \begin{list}%
   {\makebox[1cm][c]{\choicebox}\hspace{3mm}\makebox[1cm][c]{\choicebox}}%
   {\setlength{\labelwidth}{2.31 cm}\setlength{\labelsep}{3mm}
    \setlength{\leftmargin}{2.61 cm}\setlength{\listparindent}{0pt}
    \setlength{\itemindent}{0pt}}%
  }{\end{list}}

\newcounter{theaufgabe}\setcounter{theaufgabe}{1}
\newenvironment{aufgabe}[1]%
  {\bigskip\par\noindent\begin{nopagebreak}
   \textsf{\textbf{\arabic{theaufgabe}.\thinspace Exercise}}\quad
      \textsf{\textit{#1}}\\*[1ex]%
\stepcounter{theaufgabe}\hspace{2ex}\end{nopagebreak}}
  {\par\pagebreak[2]}

\renewcommand{\labelenumi}{\alph{enumi})}
\renewcommand{\labelenumii}{\arabic{enumii})}
\newcommand{\abgabe}{\marginpar{$\Box$}}
\reversemarginpar
\setlength{\parindent}{0.0cm}
\setcounter{secnumdepth}{0}

%% One-page fit: spacing only, no content changes ---------------------------
\usepackage{enumitem}
\setlist[itemize]{topsep=1pt,partopsep=0pt,itemsep=1pt,parsep=0pt,leftmargin=1.4em}
\setlist[enumerate]{topsep=1pt,partopsep=0pt,itemsep=1pt,parsep=0pt,leftmargin=1.6em}
\linespread{0.94}
\renewcommand{\arraystretch}{0.95}
%% --------------------------------------------------------------------------

\begin{document}
\selectlanguage{english}
\pagestyle{plain}

%=== Header ====================================================================
\thispagestyle{empty}
\noindent
\begin{minipage}[b][1.4cm]{1.0\textwidth}
\begin{center}
\begin{bf}
\begin{large} Structural Bioinformatics SS 2026\end{large} \\[0.15cm]
\begin{Large} Challenge: Protein Bindingpocket Identification\end{Large} \\[0.15cm]
\end{bf}
\begin{large}
Amr Abdullah, \textless k12340330 \textgreater\\
\end{large}
\end{center}
\end{minipage}

\noindent \rule[0.8em]{\textwidth}{0.12mm}\\[-0.5em]
%=====================================================================

\textbf{Task \& metric.} For every residue of a protein, decide whether it is
part of a binding pocket. The score is a flat (micro-averaged)
Jaccard index / IoU: all \texttt{(protein, residue)} pairs of the test set are
pooled into a single confusion matrix.

\textbf{Input data \& pre-processing.} 16{,}379 training proteins -with a list
of binding residues each- and 3{,}064 test proteins, all as PDB files. We read
the \texttt{ATOM} records directly without external structure library. Every
residue that has a C$\alpha$ atom becomes one example; about 11\% are binding
(on average $\approx$50 binding residues per protein). The train/validation
split is done \emph{by protein} (10\% held out), so no residue of a validation
protein is ever seen during training.

\textbf{Features (66 per residue).} Computed only from the C$\alpha$ coordinates
and residue chemistry:
\begin{itemize}
  \item \textbf{Identity/chemistry (30):} amino-acid one-hot (21)
  \item \textbf{Local geometry (22):} neighbor counts at 6/8/10/12/15\,\AA,
        NN distance statistics, half-sphere exposure, local
        anisotropy, side-chain extent, B-factor statistics, and
        \emph{pocket concavity / open-space} scores plus depth and
        solvent-accessibility proxies.
  \item \textbf{Pocket chemistry (14):} mean of 7 chemical properties over the
        C$\alpha$ neighbours within 8\,\AA{} and within 12\,\AA.
\end{itemize}

\textbf{Models:}
\begin{enumerate}
  \item \textbf{Gradient boosting}
        \texttt{HistGradientBoostingClassifier} (scikit-learn) on the 66
        features. Hyperparameters: 800 max iterations, learning rate 0.05, 63
        max leaves, min 80 samples/leaf, L2 = 1.0, class-balanced loss, early
        stopping on an internal 10\% split.
  \item \textbf{Off-the-shelf tools} \emph{P2Rank} (pre-trained random forest;
        best strategy: union of the top-2 predicted pockets) and \emph{Fpocket}. DeepSite/DeepSurf were surveyed but not practical to run at this scale.
  \item \textbf{Equivariant GNN } A 5-layer E(n)-equivariant GNN
        over a $k$-NN graph of C$\alpha$ atoms (hidden 96, $k$=16, 28-dim
        chemistry-only node features, $\approx$340k parameters, BCE loss with
        \texttt{pos\_weight}=8, Adam lr $10^{-3}$).
\end{enumerate}

\textbf{Final approach.} The gradient boosting model. At prediction time each
residue gets a probability; the global threshold (0.725) is the one that
maximizes the held-out flat Jaccard.

\textbf{Post-processing.} Two cheap operations on the C$\alpha$ graph, tuned on
the held-out proteins: (1) \emph{dilation}: promote an uncertain residue
($p\!\ge\!0.6$) to positive if it lies within 6\,\AA{} of a confident one
($p\!\ge\!0.8$), which fills pocket rims; (2) \emph{pruning}: drop connected
components of fewer than 6 predicted residues, removing isolated false
positives. This raises the held-out flat Jaccard from 0.3678 to 0.3706
(improvement confirmed on 99\% of bootstrap resamples).

\textbf{Results.}
\begin{center}
\begin{tabular}{lcc}
\toprule
Approach & Local Jaccard & Public LB \\
\midrule
\textbf{Gradient boosting + dilation and post-processing (final)} & \textbf{0.371} & \textbf{0.331} \\
Gradient boosting without V1 & \textbf{0.370} & \textbf{0.326} \\
P2Rank top-2 pockets, pre-trained          & 0.282 & 0.268 \\
Equivariant GNN              & 0.172 & 0.163 \\
\bottomrule
\end{tabular}
\end{center}

\textbf{Code \& reproducibility.} Main file (self-contained, runs end-to-end):
\texttt{01\_gbm\_handcrafted/binding\_site\_gbm.ipynb}; the same logic is also in
the command-line scripts under \texttt{01\_gbm\_handcrafted/src/}. The other two
approaches are in \texttt{02\_offshelf\_tools/} and \texttt{03\_egnn/}.

\textbf{Sources} scikit-learn
(\texttt{HistGradientBoostingClassifier}, \texttt{jaccard\_score}); P2Rank
(Kriv\'ak \& Hoksza, \emph{J.\ Cheminform.}, 2018); Fpocket (Le Guilloux et al.,
\emph{BMC Bioinformatics}, 2009); EGNN architecture (Satorras, Hoogeboom \&
Welling, \emph{E(n) Equivariant Graph Neural Networks}, ICML 2021 --- our own
implementation). All feature engineering and post-processing are original.

\end{document}
